\section{第06章 ソフトウェアエンジニアリング運用}
\begin{pointbox}{この資料の主張}
  ソフトウェアエンジニアリング運用は、完成したソフトウェアを「本番に置く作業」だけではない。運用は、配備、設定、監視、障害対応、変更管理、容量管理、バックアップ、災害復旧、サービス水準、セキュリティ、自動化を含む一連の工学活動である。現代の運用では、開発と運用を分離せず、開発運用連携（DevOps）、インフラストラクチャのコード化（Infrastructure as Code, IaC）、サイト信頼性エンジニアリング（Site Reliability Engineering, SRE）などを使い、ソフトウェアを安全かつ継続的に改善することが重要である。
  \end{pointbox}

\begin{pointbox}{この資料の読み方}
  専門用語は原則として日本語で示し、必要に応じて英語を括弧内に併記する。英語名は、元資料、標準、ツール名を調べる手がかりとして残す。初学者が前提知識なしに読めるように、各節では「何のための概念か」「実務では何を見るべきか」「失敗すると何が起こるか」を重視する。文体は、だである調で統一する。
  \end{pointbox}
